\FloatBarrier
\newpage
\section{假设检验}
\label{sec:prob-testing}

估计直接寻找参数的可能位置。另一类任务先提出一个模型假设，再比较观测结果与该假设下的随机波动。如果预先选定的偏离在这个模型下很少发生，就据此作出拒绝假设的判断。

\subsection{假设与随机波动}

设独立抛一枚硬币10次，\(K\)为正面次数，未知正面概率为\(p\)。为研究硬币是否偏向正面，提出
\[
 H_0:p=\frac12,\qquad H_1:p>\frac12.
\]
\(H_0\)称为{\heiti 原假设}，\(H_1\)为所考察的{\heiti 备择假设}。在\(H_0\)下，\(K\sim B(10,1/2)\)，因此可以计算各种观测的概率。

偏向正面的证据由较大的\(K\)提供。开始观察前，应确定用哪个统计量、按哪个方向偏离以及怎样作出判断。观察后再根据结果选择方向，会改变原来概率计算所对应的程序。

假设检验始终需要保留随机波动的可能性。即使\(H_0\)成立，也可能出现很大的\(K\)；下面通过拒绝规则控制这种错误发生的概率。

\subsection{拒绝域}

\begin{definition}
若统计量落入预先指定的集合\(R\)就拒绝\(H_0\)，称\(R\)为{\heiti 拒绝域}。若\(H_0\)成立时拒绝的概率不超过\(\alpha\)，称此规则的显著性水平不超过\(\alpha\)。
\end{definition}

在前面的硬币模型中，选择\(R=\{K\ge9\}\)，则
\[
 P_{1/2}(K\ge9)=\frac{\binom{10}9+\binom{10}{10}}{2^{10}}
 =\frac{11}{1024}\approx0.0107.
\]
若扩大为\(K\ge8\)，概率变为\(56/1024\approx0.0547\)。因此要求水平不超过\(0.05\)时，这类整数阈值规则可选择\(K\ge9\)，实际错误概率小于给定上限。

观测到\(K=9\)，便按既定规则拒绝公平假设，认为数据提供了偏向正面的证据。观测到\(K=7\)时不拒绝，表示当前数据在这一规则下尚不足以支持拒绝；这不等于证明\(p=1/2\)。

若改问硬币是否公平，备择为\(H_1:p\ne1/2\)，正面次数过大或过小都构成偏离。取双侧拒绝域\(K\le1\)或\(K\ge9\)，原假设下两端概率相加为\(22/1024\approx0.0215\)。

\begin{example}
某工艺原不合格率为\(0.1\)。为判断不合格率是否升高，独立抽检20件，预先规定发现至少5件不合格就拒绝原假设。写出原假设、备择、统计量和拒绝域，计算原假设成立时错误拒绝的概率。
\end{example}

\subsection{列联表}

两个分类变量的观测频数可以整理为\(2\times2\)列联表：
\[
\begin{array}{c|cc|c}
 & B&B^c&\text{合计}\\\hline
 A&a&b&a+b\\
 A^c&c&d&c+d\\\hline
 \text{合计}&a+c&b+d&N
\end{array}
\]
现在考察总体中两分类变量是否独立。原假设为\(P(A\cap B)=P(A)P(B)\)。用样本边际比例估计两类概率，独立假设下第一格的期望频数为\((a+b)(a+c)/N\)，其余格同理由行合计乘列合计再除以\(N\)得到。

记观测频数为\(O_{ij}\)，拟合的期望频数为\(E_{ij}\)。{\heiti Pearson统计量}定义为
\begin{equation}\label{eq:prob-pearson}
 \chi^2=\sum_{i,j}\frac{(O_{ij}-E_{ij})^2}{E_{ij}}
 =\frac{N(ad-bc)^2}{(a+b)(c+d)(a+c)(b+d)}.
\end{equation}
用样本边际拟合后，观测与期望的各行、各列之差都合计为0。因此只要知道第一格的差\(\delta\)，其余三格也随之确定：
\[
 \delta=a-\frac{(a+b)(a+c)}N=\frac{ad-bc}N,
 \qquad (O_{ij}-E_{ij})=
 \begin{pmatrix}\delta&-\delta\\-\delta&\delta\end{pmatrix}.
\]
把四个平方差代回统计量，得到\(\chi^2=\delta^2\sum_{i,j}1/E_{ij}\)，通分便得式\eqref{eq:prob-pearson}。\(ad=bc\)时，四格都恰好等于拟合期望；偏离越大，统计量越大。

标准正态变量\(Z\)的平方\(Z^2\)所服从的分布，称为自由度为1的\(\chi^2\)分布。在独立观测、各格期望频数足够大的抽样模型中，原假设下统计量\(\chi^2\)近似服从此分布。本章引用这一近似结论及其临界概率，使用时要求各格拟合期望频数至少为5。常用表为
\[
\begin{array}{c|ccc}
 \alpha&0.05&0.01&0.001\\\hline
 \chi^2_\alpha&3.841&6.635&10.828
\end{array}
\]
其含义是原假设下\(P(\chi^2\ge\chi^2_\alpha)\approx\alpha\)。据此在\(\chi^2\ge\chi^2_\alpha\)时拒绝独立假设，结论的水平由这个假设下的尾概率解释。

样本量也影响这一数值。若另一次调查的四格频数恰为原表的\(m\)倍，各组比例保持不变，而式\eqref{eq:prob-pearson}给出新统计量为\(m\chi^2\)。同样的比例差异，在较大的独立样本中可以提供更强的拒绝证据；统计量本身不能直接充当关联强弱的尺度。

\begin{example}[据2025年全国I卷改编]\label{ex:prob-medical-table}
从做过超声波检查的人群中随机调查1000人，得到下表。
\[
\begin{array}{c|rr|r}
 &\text{正常}&\text{不正常}&\text{合计}\\\hline
 \text{患该疾病}&20&180&200\\
 \text{未患该疾病}&780&20&800\\\hline
 \text{合计}&800&200&1000
\end{array}
\]
用频率估计检查结果不正常者患该疾病的概率。再根据\(\alpha=0.001\)的独立性检验，分析检查结果与是否患该疾病的关系，说明判断对应的人群范围。
\end{example}

\begin{exercise}
\mbox{}\par\nobreak

（1）在10次抛币模型中，分别以\(K\ge8\)、\(K\ge9\)为拒绝域，比较原假设下的错误拒绝概率。

（2）两种教学安排各观察50人，甲组达标40人、未达标10人，乙组达标30人、未达标20人。写出列联表及期望频数，在\(\alpha=0.05\)水平作独立性检验。

（3）说明为什么“在1\%水平拒绝独立假设”不能解释为“两个变量独立的概率只有1\%”。

（4）某调查仅记录自愿填写网络问卷的人。即使计算得到很大的Pearson统计量，把结论用于目标总体以前，还需检查什么数据获取条件？

（5）把第（2）题四格频数全部加倍，比较各组达标比例及Pearson统计量。在\(\alpha=0.01\)水平，原表与新表的检验结论是否相同？

（6）由\(Z^2\)的分布及正态对称性，说明自由度为1的\(\chi^2\)分布满足\(P(Z^2\ge c)=2[1-\Phi(\sqrt c)]\)。用正态概率表说明前两个临界值的来源。

（7）两组各\(n\)人，甲、乙达标比例为\(60\%\)、\(40\%\)，\(n\)为5的倍数。求\(\chi^2\)关于\(n\)的表达式，并求在\(0.05\)水平拒绝独立假设的最小\(n\)，核对期望频数条件。
\end{exercise}

\FloatBarrier
\newpage
\subsection{检验误差*}

原假设成立时错误拒绝，称为{\heiti 第一类错误}；某个备择模型成立时未能拒绝原假设，称为{\heiti 第二类错误}。在给定备择参数\(p\)下，拒绝概率称为该处的{\heiti 检验效能}，等于1减去第二类错误概率。

对\(K\ge9\)的硬币规则，若实际\(p=0.8\)，效能为
\[
 P_{0.8}(K\ge9)=\binom{10}9(0.8)^9(0.2)+(0.8)^{10},
\]
第二类错误概率则为其补数。扩大拒绝域会提高发现这一偏离的机会，同时也增加原假设成立时的错误拒绝。样本量和所考察备择的差异，共同影响两种错误的取舍。

在规则\(K\ge c\)下，把拒绝概率写成参数的函数，便得到
\[
 \pi_{n,c}(p)=\sum_{k=c}^n\binom nkp^k(1-p)^{n-k}.
\]
同一个函数在\(p=1/2\)处给出错误拒绝概率，在\(p>1/2\)处给出发现相应偏离的概率。例如比较以下两条事先选定的规则：
\[
 \begin{array}{c|cc}
 \text{拒绝规则}&p=1/2\text{时拒绝的概率}&p=0.8\text{时拒绝的概率}\\\hline
 n=10,\ K\ge9&0.0107&0.3758\\
 n=20,\ K\ge15&0.0207&0.8042
 \end{array}
\]
两者都把原假设下的错误概率控制在\(0.05\)以内，后一规则发现\(p=0.8\)这一偏离的机会更大。小的显著性水平只约束原假设下的表现；能否发现所关心的偏离，还要考察备择下的抽样分布。

对预先确定的单侧大值偏离，若观测到\(K=k\)，定义{\heiti p值}为\(P_{H_0}(K\ge k)\)。它计算原假设下至少同样偏大的结果的概率。在10次抛币模型中，观察到9次正面的p值为\(11/1024\)，观察到8次则为\(56/1024\)。按给定水平比较p值与\(\alpha\)，与对应的尾部拒绝域规则一致。

\subsection{条件检验*}

继续研究列联表。固定两行的对象数，并在独立假设下把第二类变量视为共同成功概率\(p\)的独立0、1记录。给定全部成功数后，任何具有该成功总数的标记方式都有相同概率因子\(p^M(1-p)^{N-M}\)，故条件下这些标记方式等可能。

记总数为\(N\)、成功总数为\(M\)，第一行有\(n\)个对象，其成功数\(X\)的条件分布为
\begin{equation}\label{eq:prob-fisher-conditional}
 \begin{aligned}
 P(X=k\mid\text{边际合计})
 &=\frac{\binom nk\binom{N-n}{M-k}}{\binom NM}\\
 &=\frac{\binom Mk\binom{N-M}{n-k}}{\binom Nn}.
 \end{aligned}
\end{equation}
未知的共同成功概率从条件分布中消去。若事先研究第一行成功比例偏高的方向，就用\(X\ge x_{\rm obs}\)的条件尾概率作判断。

例如两组各6个对象，总成功数6个，第一组成功5个，单侧尾概率为
\[
 \frac{\binom65\binom61+\binom66\binom60}{\binom{12}6}
 =\frac{37}{924}\approx0.0400.
\]
它直接由有限条件模型得到，不需要大频数近似。对例\ref{ex:prob-medical-table}，若按异常检查结果为成功，则固定边际后的\(N=1000,M=200,n=200\)，观测值为180；同一公式给出相应的条件尾概率表达。

边际固定时，列联表只有第一格\(X\)可以变化，其在独立模型下的条件期望为\(nM/N\)。将四格写成\(X,n-X,M-X,N-n-M+X\)，代入式\eqref{eq:prob-pearson}，得到
\[
 \chi^2=\frac{N^3}{n(N-n)M(N-M)}
 \left(X-\frac{nM}N\right)^2.
\]
因此两种方法都研究这一格偏离独立模型中心的程度。条件检验直接使用超几何尾概率，Pearson检验则使用前面引用的大频数近似；二者在同一张表上有共同的研究对象。

\begin{exercise}
\mbox{}\par\nobreak

（1）对硬币模型取备择\(p=0.7\)，计算\(K\ge9\)规则的效能，与\(p=0.8\)比较。

（2）两组各5个对象，总成功数4个，第一组成功4个。写出第一组成功数的条件分布，求单侧尾概率。

（3）在已知方差的正态均值模型中，证明双侧置信区间不包含\(\mu_0\)，等价于标准化样本均值落在相应的双侧尾部。由此构造检验\(H_0:\mu=\mu_0\)的规则。

（4）对\(1\le c\le n\)，通过求导与合并相邻项，证明\(\pi'_{n,c}(p)=n\binom{n-1}{c-1}p^{c-1}(1-p)^{n-c}\)。这说明成功概率升高时，单侧规则的效能怎样变化？

（5）两组各6个对象、总成功数为6时，第一组成功5个与成功1个分别对应哪个方向的偏离？若事先考察两组成功概率是否不同，采用\(X\le1\)或\(X\ge5\)的拒绝域，求其条件概率。
\end{exercise}
